

\documentclass[12pt, a4paper]{article}

% Support de l'arabe (compiler avec XeLaTeX)
\usepackage{polyglossia}
\setmainlanguage{arabic}
\setotherlanguage{french}

% Polices arabes
\newfontfamily\arabicfont[Script=Arabic,Scale=1.1]{Amiri}
\newfontfamily\arabicfontsf[Script=Arabic,Scale=1.1]{Traditional Arabic}

\usepackage{geometry}
\usepackage{setspace}
\usepackage{graphicx}

% Marges et interligne
\geometry{top=2.5cm, bottom=2.5cm, left=2.5cm, right=2.5cm}
\setstretch{1.35}

\begin{document}

\thispagestyle{empty}

% --- Titre ---
\begin{center}
    \Huge\textbf{ملخص}
\end{center}

\vspace{1cm}

% --- Corps du résumé ---
\begin{center}
\vspace{0.5cm}
\large

هذا نص تجريبي يُستخدم كعنصر نائب في هذا القسم من الوثيقة. يهدف هذا النص إلى محاكاة شكل الملخص وتوزيع الأسطر والفقرات، دون تقديم أي معلومات فعلية أو مرتبطة بموضوع محدد. يمكن استبداله لاحقاً بالمحتوى النهائي الخاص بالمشروع.\\

\vspace{0.5cm}

يتناول هذا الجزء مجموعة من الأفكار العامة المتعلقة بالموضوع المدروس، مع عرض موجز للسياق والأهداف والمنهجية المتبعة. وقد تم تنظيم المعلومات بطريقة تسمح بتقديم نظرة شاملة ومختصرة حول مختلف مراحل العمل والنتائج المتوقعة.\\

\vspace{0.5cm}

كما يمكن أن يتضمن هذا القسم وصفاً عاماً للحلول والأدوات المستخدمة، بالإضافة إلى توضيح الجوانب التقنية والمنهجية التي تم اعتمادها خلال تنفيذ المشروع. ويُترك هذا النص مؤقتاً من أجل الحفاظ على البنية والتنسيق العام للوثيقة.\\

\vspace{0.5cm}

وأخيراً، يسلط هذا القسم الضوء بشكل عام على النتائج والاستنتاجات الرئيسية، مع الإشارة إلى بعض الآفاق المستقبلية التي يمكن تطويرها في إطار الأعمال اللاحقة.\\

\vspace{1.5cm}

\textbf{الكلمات المفتاحية:} كلمة مفتاحية، كلمة أخرى، موضوع الدراسة، المنهجية، النتائج، الآفاق المستقبلية.

\end{center}

\end{document}
